%% ──────────────────────────────────────────────
%%  فصل ۱۱ – روان‌شناسی اجتماعی نجات‌طلبی
%%  فایل: chapters/ch11-social-psychology.tex
%% ──────────────────────────────────────────────
\chapter{روان‌شناسی اجتماعی نجات‌طلبی}
\label{ch:social-psychology}

%% ── چکیده‌ی فصل ──
\begin{tcolorbox}[mybox, title={چکیده‌ی فصل}]
فصل‌های پیشین این کتاب عمدتاً بر ابعاد تاریخی، حقوقی، سیاسی، و ژئوپولیتیک مداخلات خارجی و نجات‌طلبی تمرکز داشتند. اما پرسشی بنیادین همچنان بی‌پاسخ مانده است: \textbf{چرا ملت‌ها — علی‌رغم تجربه‌های مکرر شکست — باز هم به «نجات‌بخش خارجی» امید می‌بندند؟} پاسخ این پرسش را باید در لایه‌های عمیق‌تر روان‌شناسی فردی و اجتماعی جست‌وجو کرد. این فصل به بررسی مکانیزم‌های روان‌شناختی‌ای می‌پردازد که نجات‌طلبی را تولید، بازتولید، و تقویت می‌کنند: از \concept{ناامیدی آموخته‌شده}ی سلیگمن تا \concept{آرزواندیشی} و سوگیری‌های شناختی، از \concept{روان‌شناسی قربانی} تا \concept{تروماهای جمعی}، از \concept{نظریه‌ی هویت اجتماعی} تا \concept{آزمایش اطاعت} میلگرام. نشان خواهیم داد که نجات‌طلبی نه یک «انتخاب عقلانی» بلکه اغلب محصول فرآیندهای روان‌شناختی قابل پیش‌بینی — و قابل بهره‌برداری — است.
\end{tcolorbox}

%% ================================================================
\section{درآمد: چرا روان‌شناسی؟}
\label{sec:ch11-intro}
%% ================================================================

\needspace{6\baselineskip}
در فصل‌های ۴ تا ۱۰، ده‌ها مورد تاریخی مداخله‌ی خارجی را بررسی کردیم. یک الگوی تکرارشونده در همه‌ی این موارد دیده می‌شود: گروهی از مردم تحت فشار (سیاسی، اقتصادی، نظامی) به نیروی خارجی چشم می‌دوزند، از آن کمک می‌خواهند یا مداخله‌اش را خوشامد می‌گویند، اما نتیجه‌ی نهایی اغلب بدتر از وضعیت اولیه است. با این حال، این الگو بارها و بارها تکرار می‌شود.

\thinker{آلبرت اینشتین} (\lat{Albert Einstein}) — به نقل منسوب — دیوانگی را «تکرار یک کار و انتظار نتیجه‌ی متفاوت» تعریف کرده است. آیا ملت‌های نجات‌طلب «دیوانه»اند؟ قطعاً نه. اما ساختارهای شناختی و عاطفی انسان به‌گونه‌ای طراحی شده‌اند که در شرایط فشار و ناامیدی، الگوهای رفتاری خاصی را — حتی علی‌رغم شواهد عقلانی بر خلاف آن — تکرار کنند.

\thinker{دانیل کانمن} (\lat{Daniel Kahneman})، روان‌شناس برنده‌ی جایزه‌ی نوبل اقتصاد، در کتاب اثرگذار خود \textit{تفکر، سریع و کند} نشان داده است که ذهن انسان از دو سیستم تشکیل شده: \concept{سیستم ۱} (سریع، شهودی، احساسی) و \concept{سیستم ۲} (کند، تحلیلی، منطقی). در شرایط بحران و فشار، سیستم ۱ غالب می‌شود و سوگیری‌های شناختی تشدید می‌شوند.\footnote{\lr{Kahneman, Daniel. \textit{Thinking, Fast and Slow}. New York: Farrar, Straus and Giroux, 2011, pp. 19–30.}}

\begin{tcolorbox}[defbox, title={مفهوم کلیدی: روان‌شناسی سیاسی نجات‌طلبی}]
\concept{روان‌شناسی سیاسی نجات‌طلبی} به مجموعه‌ی فرآیندهای شناختی، عاطفی، و اجتماعی‌ای اطلاق می‌شود که در جوامع تحت فشار، زمینه‌ی پذیرش، دعوت، یا استقبال از مداخله‌ی خارجی را فراهم می‌آورند. این فرآیندها شامل سوگیری‌های شناختی (آرزواندیشی، تأیید، خوش‌بینی)، پویایی‌های گروهی (هویت اجتماعی، قطب‌بندی، اطاعت)، و آسیب‌های روان‌شناختی (تروما، ناامیدی آموخته‌شده، سندرم قربانی) می‌شوند.
\end{tcolorbox}

این فصل بر اساس آثار روان‌شناسان اجتماعی، شناختی، و سیاسی از جمله \thinker{مارتین سلیگمن} (\lat{Martin Seligman})، \thinker{هنری تاجفل} (\lat{Henri Tajfel}) و \thinker{جان ترنر} (\lat{John Turner})، \thinker{استنلی میلگرام} (\lat{Stanley Milgram})، \thinker{فیلیپ زیمباردو} (\lat{Philip Zimbardo})، \thinker{ایروینگ جنیس} (\lat{Irving Janis})، و \thinker{ولکان واموک} (\lat{Vamık Volkan}) سازمان‌دهی شده است.


%% ================================================================
\section{ناامیدی آموخته‌شده: چرا ملت‌ها دست از تلاش برمی‌دارند؟}
\label{sec:ch11-learned-helplessness}
%% ================================================================

\needspace{6\baselineskip}
\subsection{نظریه‌ی اصلی سلیگمن}
\label{subsec:ch11-seligman-theory}

در سال ۱۹۶۷، \thinker{مارتین سلیگمن} (\lat{Martin Seligman}) و \thinker{استیون مایر} (\lat{Steven Maier}) در آزمایش‌های مشهور خود بر روی سگ‌ها، پدیده‌ی \concept{ناامیدی آموخته‌شده} (\lat{Learned Helplessness}) را کشف کردند. سگ‌هایی که بارها در معرض شوک الکتریکی غیرقابل‌اجتناب قرار گرفته بودند، حتی وقتی امکان فرار فراهم شد، تلاش نکردند — آن‌ها «آموخته» بودند که هیچ کاری نمی‌توانند بکنند.\footnote{\lr{Seligman, Martin E.P. and Steven F. Maier. "Failure to Escape Traumatic Shock," \textit{Journal of Experimental Psychology}, Vol. 74, No. 1, 1967, pp. 1–9.}}

\begin{tcolorbox}[defbox, title={تعریف: ناامیدی آموخته‌شده}]
\concept{ناامیدی آموخته‌شده} حالتی است که در آن فرد (یا گروه) پس از تجربه‌ی مکرر شکست و ناکامی در کنترل محیط، باور می‌کند که هیچ اقدامی نتیجه‌بخش نیست و بنابراین از تلاش دست برمی‌دارد — حتی زمانی که شرایط عینی تغییر کرده و امکان تغییر وجود دارد. سه مؤلفه‌ی اصلی آن عبارت‌اند از:
\begin{enumerate}
\item \textbf{نقص انگیزشی:} کاهش تلاش و ابتکار عمل
\item \textbf{نقص شناختی:} ناتوانی در یادگیری اینکه پاسخ‌های جدید می‌توانند مؤثر باشند
\item \textbf{نقص عاطفی:} افسردگی، بی‌تفاوتی، و انفعال
\end{enumerate}
\end{tcolorbox}

سلیگمن بعدها نظریه‌اش را به انسان‌ها تعمیم داد و نشان داد که افسردگی بالینی ارتباط نزدیکی با ناامیدی آموخته‌شده دارد.\footnote{\lr{Seligman, Martin E.P. \textit{Helplessness: On Depression, Development, and Death}. San Francisco: W.H. Freeman, 1975.}} در ویرایش بازنگری‌شده‌ی نظریه (با \thinker{لین آبرامسون} (\lat{Lyn Abramson}) و \thinker{جان تیزدیل} (\lat{John Teasdale}))، نقش \concept{سبک اسنادی} (\lat{Attributional Style}) برجسته شد: افرادی که شکست‌ها را به عوامل \textit{درونی} (تقصیر خودم)، \textit{پایدار} (همیشه همین‌طور بوده)، و \textit{فراگیر} (همه‌جا همین‌طور است) نسبت می‌دهند، بیشتر مستعد ناامیدی آموخته‌شده هستند.\footnote{\lr{Abramson, Lyn Y., Martin E.P. Seligman, and John D. Teasdale. "Learned Helplessness in Humans: Critique and Reformulation," \textit{Journal of Abnormal Psychology}, Vol. 87, No. 1, 1978, pp. 49–74.}}

\subsection{تعمیم به سطح جمعی: ملت‌های ناامید}
\label{subsec:ch11-collective-helplessness}

اگرچه سلیگمن نظریه‌اش را در سطح فردی توسعه داد، \thinker{مارتین بار} (\lat{Ignacio Martín-Baró})، روان‌شناس اجتماعی اهل السالوادور (که خود قربانی ترور شد)، مفهوم \concept{تروما-زدگی روانی-اجتماعی} (\lat{Psychosocial Trauma}) را برای تحلیل جوامع تحت سرکوب ارائه کرد. از نظر او، ملت‌هایی که سال‌ها تحت دیکتاتوری، جنگ، یا استعمار بوده‌اند، ویژگی‌های ناامیدی آموخته‌شده را در سطح جمعی نشان می‌دهند.\footnote{\lr{Martín-Baró, Ignacio. \textit{Writings for a Liberation Psychology}. Edited by Adrianne Aron and Shawn Corne. Cambridge, MA: Harvard University Press, 1994, pp. 108–135.}}

\begin{tcolorbox}[casebox, title={مطالعه‌ی موردی: ناامیدی آموخته‌شده در عراق}]
عراق پس از دهه‌ها دیکتاتوری بعثی، سه جنگ ویرانگر (ایران-عراق، کویت، ۲۰۰۳)، و دوازده سال تحریم‌های فلج‌کننده، نمونه‌ای بارز از ناامیدی آموخته‌شده‌ی جمعی بود:

\textbf{۱. نقص انگیزشی:} بخش قابل‌توجهی از مردم عراق — به‌ویژه پس از سرکوب انتفاضه‌ی ۱۹۹۱ — باور خود به امکان تغییر از درون را از دست دادند. همین ناامیدی زمینه‌ساز استقبال بخشی از مردم (به‌ویژه شیعیان و کردها) از اشغال آمریکا شد.

\textbf{۲. نقص شناختی:} ناتوانی در تصور جایگزینی میان دیکتاتوری و اشغال — گویی تنها دو گزینه وجود داشت.

\textbf{۳. نقص عاطفی:} بی‌تفاوتی سیاسی گسترده و پناه بردن به هویت‌های فرقه‌ای (شیعه، سنی، کرد) به‌جای هویت ملی.

\thinker{کانعان مکیّه} (\lat{Kanan Makiya})، نویسنده‌ی عراقی‌تبار، در کتاب \textit{جمهوری ترس} نشان داده است که چگونه دستگاه سرکوب بعث، «فرهنگ ترس» جمعی‌ای ایجاد کرد که ابتکار عمل و مقاومت مستقل را فلج ساخت.\footnote{\lr{Makiya, Kanan. \textit{Republic of Fear: The Politics of Modern Iraq}. Updated ed., Berkeley: University of California Press, 1998.}}
\end{tcolorbox}

\subsection{ناامیدی آموخته‌شده و پذیرش مداخله}
\label{subsec:ch11-helplessness-intervention}

رابطه‌ی میان ناامیدی آموخته‌شده و نجات‌طلبی مستقیم و روشن است:

\begin{tcolorbox}[wavebox, title={مکانیزم: از ناامیدی تا نجات‌طلبی}]
\begin{enumerate}
\item ملت تحت سرکوب تجربه‌ی مکرر شکست در تغییر از درون را انباشت می‌کند
\item باور به «عاملیت» (\lat{agency}) خود کاهش می‌یابد: «ما نمی‌توانیم»
\item \concept{مرجع کنترل بیرونی} (\lat{External Locus of Control}) غالب می‌شود: «فقط بیرونی‌ها می‌توانند»
\item نیروی خارجی به‌عنوان تنها راه‌حل تصور می‌شود
\item هزینه‌ها و خطرات مداخله‌ی خارجی نادیده گرفته یا کم‌اهمیت تلقی می‌شوند
\item مداخله صورت می‌گیرد اما نتیجه‌ی مطلوب نمی‌دهد → ناامیدی عمیق‌تر → چرخه تکرار
\end{enumerate}
\end{tcolorbox}

\thinker{جیمز اسکات} (\lat{James C. Scott}) در کتاب \textit{سلاح‌های ضعیفان} استدلال کرده است که حتی در شرایط ناامیدی ظاهری، اَشکال پنهان مقاومت (کارشکنی، شایعه‌سازی، فرار مالیاتی) ادامه دارد. اما \concept{مقاومت پنهان} با \concept{مقاومت سازمان‌یافته} متفاوت است: اولی ناامیدی را بازتولید می‌کند، دومی آن را می‌شکند.\footnote{\lr{Scott, James C. \textit{Weapons of the Weak: Everyday Forms of Peasant Resistance}. New Haven: Yale University Press, 1985.}}


%% ================================================================
\section{آرزواندیشی و سوگیری‌های شناختی}
\label{sec:ch11-wishful-thinking}
%% ================================================================

\needspace{6\baselineskip}
\subsection{تعریف و مکانیزم آرزواندیشی}
\label{subsec:ch11-wishful-definition}

\concept{آرزواندیشی} (\lat{Wishful Thinking}) — مفهوم محوری عنوان این کتاب — به گرایش شناختی‌ای اشاره دارد که در آن فرد (یا گروه) باورهایش را نه بر اساس شواهد عینی، بلکه بر اساس آنچه \textit{دوست دارد} واقعیت داشته باشد، شکل می‌دهد.\footnote{\lr{Bastardi, Anthony, Eric Luis Uhlmann, and Lee Ross. "Wishful Thinking: Belief, Desire, and the Motivated Evaluation of Scientific Evidence," \textit{Psychological Science}, Vol. 22, No. 6, 2011, pp. 731–732.}}

\begin{tcolorbox}[defbox, title={تعریف: آرزواندیشی}]
\concept{آرزواندیشی} (\lat{Wishful Thinking}) شکل‌گیری یا حفظ باورها بر اساس آنچه لذت‌بخش یا دلپذیر تصور می‌شود، نه بر اساس شواهد، عقلانیت، یا واقعیت. در زمینه‌ی سیاسی، آرزواندیشی به معنای باور به اینکه مداخله‌ی خارجی نتیجه‌ی مطلوب خواهد داشت — بدون توجه کافی به شواهد تاریخی مبنی بر عکس آن — است.

در مقابل، \concept{خام‌اندیشی} (\lat{Naivety}) به فقدان دانش و تجربه‌ی کافی برای ارزیابی واقع‌بینانه اشاره دارد — تفاوت مهمی که در عنوان کتاب بازتاب یافته است.
\end{tcolorbox}

\thinker{زیوا کوندا} (\lat{Ziva Kunda})، روان‌شناس اجتماعی، در مقاله‌ی کلاسیک خود درباره‌ی \concept{استدلال انگیزه‌مند} (\lat{Motivated Reasoning}) نشان داده است که انسان‌ها نه‌تنها اطلاعات هماهنگ با باورهای‌شان را ترجیح می‌دهند، بلکه فعالانه اطلاعات ناسازگار را تحریف یا نادیده می‌گیرند.\footnote{\lr{Kunda, Ziva. "The Case for Motivated Reasoning," \textit{Psychological Bulletin}, Vol. 108, No. 3, 1990, pp. 480–498.}}

\subsection{سوگیری‌های شناختی مرتبط با نجات‌طلبی}
\label{subsec:ch11-biases}

آرزواندیشی تنها یکی از سوگیری‌های شناختی‌ای است که نجات‌طلبی را تغذیه می‌کنند. جدول زیر مهم‌ترین سوگیری‌های مرتبط را خلاصه می‌کند:

\begin{table}[htbp]
\centering
\caption{سوگیری‌های شناختی مرتبط با نجات‌طلبی}
\label{tab:ch11-biases}
\begin{adjustbox}{max width=\linewidth}
\begin{tabularx}{\linewidth}{
    >{\raggedleft\arraybackslash}p{2.8cm}
    >{\raggedleft\arraybackslash}p{3.5cm}
    >{\raggedleft\arraybackslash}X
    >{\raggedleft\arraybackslash}p{3cm}
}
\toprule
\rowcolor{PrimaryDeep}
\textcolor{white}{\textbf{سوگیری}} &
\textcolor{white}{\textbf{تعریف}} &
\textcolor{white}{\textbf{مصداق در نجات‌طلبی}} &
\textcolor{white}{\textbf{مرجع}} \\
\midrule
آرزواندیشی & باور بر اساس آرزو، نه شواهد & «آمریکا ما را نجات خواهد داد» &  \lr{Bastardi et al., 2011} \\
\rowcolor{NeutralLight}
سوگیری تأیید & جستجوی اطلاعات هماهنگ با باور & فقط اخبار مثبت درباره‌ی مداخله‌گر & \lr{Nickerson, 1998} \\
سوگیری خوش‌بینی & دست‌کم گرفتن ریسک‌ها & «مداخله کوتاه و بی‌هزینه خواهد بود» & \lr{Sharot, 2011} \\
\rowcolor{NeutralLight}
اثر هاله‌ای & تعمیم یک ویژگی مثبت به کل & «آمریکا قوی است، پس عادل هم هست» & \lr{Thorndike, 1920} \\
تفکر سیاه‌وسفید & دوقطبی‌سازی افراطی & «یا رژیم فعلی یا نجات خارجی — راه سومی نیست» & \lr{Beck, 1976} \\
\rowcolor{NeutralLight}
خطای برنامه‌ریزی & کم‌برآورد زمان و هزینه & «مداخله ماه‌ها طول می‌کشد، نه سال‌ها» & \lr{Kahneman \& Tversky, 1979} \\
سوگیری بازمانده & توجه فقط به موارد موفق & «آلمان و ژاپن پس از جنگ جهانی دوم موفق شدند» & \lr{Taleb, 2007} \\
\rowcolor{NeutralLight}
اثر عقب‌نگری & «من از اول می‌دانستم» & توجیه شکست مداخله پس از وقوع & \lr{Fischhoff, 1975} \\
مغالطه‌ی هزینه‌ی غرق‌شده & ادامه‌ی سرمایه‌گذاری به دلیل هزینه‌های قبلی & «چون تا حالا هزینه کرده‌ایم، باید ادامه دهیم» & \lr{Arkes \& Blumer, 1985} \\
\bottomrule
\end{tabularx}
\end{adjustbox}
\end{table}

\subsubsection{سوگیری تأیید و حباب اطلاعاتی}
\label{subsubsec:ch11-confirmation}

\concept{سوگیری تأیید} (\lat{Confirmation Bias}) — گرایش به جستجو، تفسیر، و یادآوری اطلاعاتی که باورهای موجود را تأیید می‌کنند — یکی از قوی‌ترین و پردامنه‌ترین سوگیری‌های شناختی است.\footnote{\lr{Nickerson, Raymond S. "Confirmation Bias: A Ubiquitous Phenomenon in Many Guises," \textit{Review of General Psychology}, Vol. 2, No. 2, 1998, pp. 175–220.}}

در زمینه‌ی نجات‌طلبی، سوگیری تأیید به دو صورت عمل می‌کند:

\begin{itemize}
\item \textbf{در سطح ملت هدف:} گروه‌هایی که به مداخله‌ی خارجی امید بسته‌اند، فقط اخبار و تحلیل‌هایی را می‌خوانند و می‌پذیرند که مداخله را مثبت نشان می‌دهد. شبکه‌های اجتماعی این «حباب اطلاعاتی» (\lat{Filter Bubble}) را تشدید می‌کنند.\footnote{\lr{Pariser, Eli. \textit{The Filter Bubble: What the Internet Is Hiding from You}. New York: Penguin Press, 2011.}}

\item \textbf{در سطح مداخله‌گر:} همان‌گونه که در فصل~\ref{ch:contemporary} دیدیم، دولت آمریکا در مورد سلاح‌های کشتار جمعی عراق دچار سوگیری تأیید سازمانی شد: اطلاعات تأییدکننده پذیرفته و اطلاعات ناقض رد شدند.\footnote{\lr{Jervis, Robert. \textit{Why Intelligence Fails: Lessons from the Iranian Revolution and the Iraq War}. Ithaca: Cornell University Press, 2010, pp. 123–152.}}
\end{itemize}

\subsubsection{سوگیری خوش‌بینی و خطای برنامه‌ریزی}
\label{subsubsec:ch11-optimism}

\thinker{تالی شاروت} (\lat{Tali Sharot}) در پژوهش‌های خود نشان داده است که مغز انسان به‌طور پیش‌فرض \concept{خوش‌بین} است: احتمال رویدادهای مثبت را بیش‌برآورد و احتمال رویدادهای منفی را کم‌برآورد می‌کند.\footnote{\lr{Sharot, Tali. \textit{The Optimism Bias: A Tour of the Irrationally Positive Brain}. New York: Pantheon Books, 2011.}} این سوگیری در سطح جمعی نیز عمل می‌کند:

\begin{tcolorbox}[quotebox, title={نقل‌قول تاریخی}]
«پسران ما تا کریسمس به خانه برمی‌گردند.»\\
— عبارت تکرارشده در آغاز تقریباً هر جنگی در تاریخ مدرن — از جنگ جهانی اول (۱۹۱۴) تا افغانستان (۲۰۰۱) و عراق (۲۰۰۳).\footnote{\lr{Kahneman, Daniel and Jonathan Renshon. "Why Hawks Win," \textit{Foreign Policy}, No. 158, 2007, pp. 34–38.}}
\end{tcolorbox}

\thinker{کانمن} و \thinker{رنشون} (\lat{Jonathan Renshon}) در مقاله‌ی مهم «چرا بازها برنده می‌شوند» نشان داده‌اند که سوگیری‌های شناختی به‌طور سیستماتیک به نفع \concept{جنگ‌طلبان} (\lat{hawks}) و به ضرر \concept{صلح‌طلبان} (\lat{doves}) عمل می‌کنند: خوش‌بینی نسبت به نتایج جنگ، بیش‌برآورد قدرت خودی، کم‌برآورد مقاومت دشمن، و خطای برنامه‌ریزی.\footnote{\lr{Kahneman, Daniel and Jonathan Renshon. "Why Hawks Win," op. cit.}}

\subsubsection{سوگیری بازمانده: «آلمان و ژاپن موفق شدند، پس...»}
\label{subsubsec:ch11-survivorship}

یکی از رایج‌ترین استدلال‌ها در دفاع از مداخله‌ی خارجی و دولت‌سازی، اشاره به نمونه‌های «موفق» آلمان و ژاپن پس از جنگ جهانی دوم است. اما این استدلال مبتلا به \concept{سوگیری بازمانده} (\lat{Survivorship Bias}) است: فقط به موارد موفق توجه می‌شود و ده‌ها مورد شکست‌خورده (ویتنام، سومالی، عراق، لیبی، افغانستان) نادیده گرفته می‌شوند.

\thinker{نسیم نیکلاس طالب} (\lat{Nassim Nicholas Taleb}) در کتاب \textit{قوی سیاه} توضیح داده است که سوگیری بازمانده یکی از خطرناک‌ترین خطاهای شناختی در تصمیم‌گیری سیاسی است: «ما فقط برندگان را می‌بینیم و از آن‌ها درس می‌گیریم، در حالی که اکثریت بازندگان در گورستان تاریخ مدفون‌اند.»\footnote{\lr{Taleb, Nassim Nicholas. \textit{The Black Swan: The Impact of the Highly Improbable}. 2nd ed., New York: Random House, 2010, pp. 101–107.}}

\begin{tcolorbox}[wavebox, title={تحلیل: چرا آلمان و ژاپن استثنا بودند، نه قاعده}]
شرایط موفقیت دولت‌سازی در آلمان و ژاپن پس از ۱۹۴۵ شامل عواملی بود که در هیچ‌یک از مداخلات معاصر وجود نداشت:
\begin{enumerate}
\item \textbf{تسلیم بی‌قیدوشرط:} جنگ با شکست کامل و بی‌ابهام پایان یافت — بر خلاف عراق و افغانستان
\item \textbf{سنت دولت مدرن:} هر دو کشور پیش از جنگ، دولت‌های مدرن، بوروکراسی کارآمد، و نظام آموزشی پیشرفته داشتند
\item \textbf{همگنی نسبی:} نه تنش‌های فرقه‌ای عراق را داشتند نه تنوع قومی افغانستان را
\item \textbf{تهدید خارجی مشترک:} جنگ سرد و تهدید شوروی، اشغالگران و مردم بومی را متحد کرد
\item \textbf{سرمایه‌گذاری عظیم:} طرح مارشال — ۱۳ میلیارد دلار (معادل حدود ۱۵۰ میلیارد دلار امروز)\footnote{\lr{Behrman, Greg. \textit{The Most Noble Adventure: The Marshall Plan and How America Helped Rebuild Europe}. New York: Free Press, 2007.}}
\item \textbf{تعهد بلندمدت:} حضور نظامی آمریکا هنوز هم ادامه دارد — هفت دهه بعد
\end{enumerate}
\end{tcolorbox}


%% ================================================================
\section{روان‌شناسی قربانی و هویت قربانی‌گری}
\label{sec:ch11-victim-psychology}
%% ================================================================

\needspace{6\baselineskip}
\subsection{هویت قربانی جمعی}
\label{subsec:ch11-collective-victimhood}

\concept{هویت قربانی جمعی} (\lat{Collective Victimhood}) به باور مشترک اعضای یک گروه مبنی بر اینکه گروهشان به‌طور ناعادلانه مورد آزار، ستم، یا تبعیض قرار گرفته اشاره دارد. \thinker{دانیل بار-تال} (\lat{Daniel Bar-Tal})، روان‌شناس اجتماعی اسرائیلی، این مفهوم را به‌طور گسترده مطالعه کرده است.\footnote{\lr{Bar-Tal, Daniel, et al. "A Sense of Self-Perceived Collective Victimhood in Intractable Conflicts," \textit{International Review of the Red Cross}, Vol. 91, No. 874, 2009, pp. 229–258.}}

بار-تال تأکید می‌کند که هویت قربانی دو عملکرد متناقض دارد:

\begin{tcolorbox}[wavebox, title={دو چهره‌ی هویت قربانی}]
\textbf{عملکرد مثبت:}
\begin{itemize}
\item انسجام گروهی و همبستگی اجتماعی
\item بسیج عاطفی برای مقاومت
\item مشروعیت‌بخشی به مبارزه‌ی حق‌طلبانه
\end{itemize}

\textbf{عملکرد منفی:}
\begin{itemize}
\item \concept{مصونیت اخلاقی ادراکی}: «چون قربانی‌ایم، هر کاری حق ماست»
\item رد مسئولیت: «تقصیر ما نیست — همه‌اش تقصیر دیگران است»
\item \concept{رقابت قربانی‌گری}: «رنج ما از رنج دیگران بیشتر است»
\item آمادگی برای پذیرش «نجات‌بخش»: «حق داریم از هر کسی کمک بخواهیم»
\end{itemize}
\end{tcolorbox}

\subsection{ولکان و مفهوم «تروما-ی انتخاب‌شده»}
\label{subsec:ch11-chosen-trauma}

\thinker{ولکان واموک} (\lat{Vamık Volkan})، روانکاو و نظریه‌پرداز ترک‌تبار آمریکایی، مفهوم بسیار مهم \concept{تروما-ی انتخاب‌شده} (\lat{Chosen Trauma}) را مطرح کرده است. منظور او رویداد تروماتیکی در تاریخ یک ملت است که به‌طور فعال در حافظه‌ی جمعی نگه‌داری، بازروایت، و بازتولید می‌شود و به بخشی از هویت گروهی تبدیل می‌گردد.\footnote{\lr{Volkan, Vamık D. \textit{Bloodlines: From Ethnic Pride to Ethnic Terrorism}. New York: Farrar, Straus and Giroux, 1997, pp. 36–49.}}

\begin{tcolorbox}[defbox, title={تعریف: تروما-ی انتخاب‌شده}]
\concept{تروما-ی انتخاب‌شده} رویداد تروماتیک تاریخی‌ای است که یک گروه قومی، ملی، یا مذهبی آن را به‌طور فعال در هویت جمعی خود ادغام کرده و از نسلی به نسل دیگر منتقل می‌کند. این تروما نه‌تنها یادآوری می‌شود، بلکه \textit{بازتجربه} می‌شود — گویی رویداد گذشته هنوز در حال وقوع است.

\textbf{نمونه‌ها:}
\begin{itemize}
\item نبرد کوزوو (۱۳۸۹) برای صرب‌ها
\item سقوط اندلس (۱۴۹۲) برای بخشی از مسلمانان
\item حمله‌ی مغول (۱۲۵۸) برای ایرانیان و عرب‌ها
\item نسل‌کشی ارمنیان (۱۹۱۵) برای ارمنی‌ها
\item هولوکاست (۱۹۳۳–۱۹۴۵) برای یهودیان
\item کودتای ۲۸ مرداد (۱۹۵۳) برای ایرانیان
\end{itemize}
\end{tcolorbox}

ولکان نشان می‌دهد که تروماهای انتخاب‌شده در زمان بحران «فعال» می‌شوند: رهبران سیاسی با بازخوانی تروما-ی تاریخی، احساسات جمعی را بسیج می‌کنند. \histref{اسلوبودان میلوشویچ}{رئیس‌جمهور صربستان} در سال ۱۹۸۹ — ششصدمین سالگرد نبرد کوزوو — سخنرانی مشهوری ایراد کرد که تروما-ی ۱۳۸۹ را برای توجیه ناسیونالیسم تهاجمی صربی به‌کار برد.\footnote{\lr{Volkan, Vamık D. "Transgenerational Transmissions and Chosen Traumas: An Aspect of Large-Group Identity," \textit{Group Analysis}, Vol. 34, No. 1, 2001, pp. 79–97.}}

\subsection{نقش تروماهای جمعی در پذیرش مداخله}
\label{subsec:ch11-trauma-intervention}

\begin{tcolorbox}[casebox, title={تحلیل: تروما و نجات‌طلبی — چهار مکانیزم}]
\textbf{مکانیزم ۱: فوریت عاطفی.}
تروماهای جمعی حس «فوریت» ایجاد می‌کنند: «باید همین الان کاری شود، حتی اگر آن کار مداخله‌ی خارجی باشد.» این فوریت مانع تحلیل عقلانی هزینه-فایده می‌شود.

\textbf{مکانیزم ۲: آرمان‌سازی نجات‌بخش.}
ملت تروماتیزه‌شده گرایش دارد نجات‌بخش بالقوه را \concept{آرمانی} (\lat{idealize}) کند — یعنی ویژگی‌های مثبت او را بزرگ‌نمایی و ویژگی‌های منفی‌اش را نادیده بگیرد. این پدیده در روانکاوی با مفهوم \concept{انتقال} (\lat{transference}) مرتبط است.\footnote{\lr{Volkan, Vamık D. \textit{Killing in the Name of Identity: A Study of Bloody Conflicts}. Charlottesville: Pitchstone Publishing, 2006, pp. 155–178.}}

\textbf{مکانیزم ۳: دوقطبی‌سازی خیر و شر.}
تروما جهان را به «خوب» و «بد» مطلق تقسیم می‌کند: رژیم فعلی = شر مطلق؛ نجات‌بخش خارجی = خیر مطلق. این \concept{شکافتن} (\lat{splitting}) — مکانیزم دفاعی بدوی در روانکاوی — مانع دیدن واقعیت پیچیده و خاکستری می‌شود.

\textbf{مکانیزم ۴: تکرار اجباری.}
\thinker{زیگموند فروید} (\lat{Sigmund Freud}) مفهوم \concept{تکرار اجباری} (\lat{Repetition Compulsion}) را مطرح کرد: فرد تروماتیزه‌شده ناخودآگاه شرایط تروماتیک را بازآفرینی می‌کند.\footnote{\lr{Freud, Sigmund. "Beyond the Pleasure Principle" (1920), in \textit{The Standard Edition of the Complete Psychological Works}, Vol. XVIII, London: Hogarth Press, 1955, pp. 1–64.}} در سطح جمعی، ملت‌های تروماتیزه ممکن است ناخودآگاه الگوهای مداخله-وابستگی-شکست را تکرار کنند.
\end{tcolorbox}


%% ================================================================
\section{نظریه‌ی هویت اجتماعی و مرز خودی/بیگانه}
\label{sec:ch11-social-identity}
%% ================================================================

\needspace{6\baselineskip}
\subsection{نظریه‌ی تاجفل و ترنر}
\label{subsec:ch11-tajfel-turner}

\thinker{هنری تاجفل} (\lat{Henri Tajfel}) و \thinker{جان ترنر} (\lat{John Turner}) در دهه‌ی ۱۹۷۰ و ۱۹۸۰ \concept{نظریه‌ی هویت اجتماعی} (\lat{Social Identity Theory — SIT}) را توسعه دادند. بر اساس این نظریه، بخش مهمی از هویت فردی از عضویت در گروه‌های اجتماعی (ملیت، قومیت، مذهب، طبقه) مشتق می‌شود. افراد گرایش دارند \concept{درون‌گروه} (\lat{in-group}) خود را ارزشمندتر از \concept{برون‌گروه} (\lat{out-group}) بدانند.\footnote{\lr{Tajfel, Henri and John C. Turner. "An Integrative Theory of Intergroup Conflict," in William G. Austin and Stephen Worchel (eds.), \textit{The Social Psychology of Intergroup Relations}. Monterey, CA: Brooks/Cole, 1979, pp. 33–47.}}

\begin{tcolorbox}[defbox, title={مفاهیم کلیدی: نظریه‌ی هویت اجتماعی}]
\begin{itemize}
\item \concept{طبقه‌بندی اجتماعی} (\lat{Social Categorization}): تقسیم جهان اجتماعی به «ما» و «آن‌ها»
\item \concept{هم‌ذات‌پنداری اجتماعی} (\lat{Social Identification}): تعریف خود بر اساس عضویت گروهی
\item \concept{مقایسه‌ی اجتماعی} (\lat{Social Comparison}): ارزیابی گروه خود در مقایسه با گروه‌های دیگر
\item \concept{تمایز مثبت} (\lat{Positive Distinctiveness}): تلاش برای حفظ تصویر مثبت از درون‌گروه
\end{itemize}
\end{tcolorbox}

\subsection{کاربرد در نجات‌طلبی: چه کسی «خودی» است؟}
\label{subsec:ch11-sit-intervention}

نظریه‌ی هویت اجتماعی کمک می‌کند بفهمیم چرا در برخی جوامع، «نجات‌بخش خارجی» پذیرفتنی‌تر از «رژیم داخلی» تلقی می‌شود. وقتی مرز خودی/بیگانه بازتعریف شود — مثلاً بر اساس مذهب به‌جای ملیت، یا بر اساس ایدئولوژی به‌جای جغرافیا — «خارجی» ممکن است «خودی‌تر» از «داخلی» به نظر برسد:

\begin{tcolorbox}[casebox, title={نمونه‌ها: بازتعریف مرز خودی/بیگانه}]
\begin{itemize}
\item \textbf{شیعیان عراق و ایران:} برای بخشی از شیعیان عراق، ایران (هم‌مذهب) «خودی‌تر» از سنی‌های عراقی بود — و بالعکس.
\item \textbf{کردها و آمریکا:} کردهای عراق، آمریکا را «دوست» و دولت مرکزی عراق را «دشمن» تلقی کردند — بازتعریف ملیت بر اساس منافع قومی.
\item \textbf{اپوزیسیون ایرانی و غرب:} بخشی از اپوزیسیون ایرانی خارج‌نشین، دولت‌های غربی را متحد طبیعی خود می‌داند — هویت‌یابی بر اساس ارزش‌ها (دموکراسی، سکولاریسم) به‌جای ملیت.
\item \textbf{صرب‌های بوسنی و صربستان:} صرب‌های بوسنی، صربستان را «مادرمیهن» و دولت بوسنی را «بیگانه» تلقی کردند.
\item \textbf{روس‌زبانان کریمه و روسیه:} روس‌زبانان کریمه (یا بخشی از آن‌ها) الحاق به روسیه را «بازگشت به خانه» تلقی کردند.
\end{itemize}
\end{tcolorbox}

این بازتعریف مرزها نه تصادفی است و نه صرفاً خودجوش: دولت‌های مداخله‌گر فعالانه در بازتعریف هویت‌ها نقش دارند. رادیوهای فارسی‌زبان، شبکه‌های تلویزیونی عربی، رسانه‌های روسی‌زبان در اوکراین — همه ابزارهایی برای بازتعریف مرز خودی/بیگانه به نفع مداخله‌گرند (موضوعی که در فصل~\ref{ch:missionary} مفصل‌تر بررسی خواهد شد).

\subsection{نظریه‌ی طبقه‌بندی خود و فراهویت}
\label{subsec:ch11-self-categorization}

\thinker{ترنر} در تکمیل نظریه‌ی هویت اجتماعی، \concept{نظریه‌ی طبقه‌بندی خود} (\lat{Self-Categorization Theory — SCT}) را ارائه کرد. بر اساس این نظریه، هویت فرد سیّال است و بسته به زمینه تغییر می‌کند: در یک زمینه فرد خود را «ایرانی» می‌داند، در زمینه‌ی دیگر «شیعه»، در زمینه‌ی سوم «کرد»، و در زمینه‌ی چهارم «دموکراسی‌خواه».\footnote{\lr{Turner, John C., et al. \textit{Rediscovering the Social Group: A Self-Categorization Theory}. Oxford: Basil Blackwell, 1987.}}

مداخله‌گران بیرونی و نجات‌طلبان داخلی هر دو تلاش می‌کنند سطح خاصی از هویت را \concept{برجسته} (\lat{salient}) سازند: مداخله‌گر می‌خواهد «ستمدیدگی» را برجسته کند (تا مداخله مشروع شود)، نجات‌طلبان می‌خواهند «هم‌ارزشی با غرب» را برجسته کنند (تا مداخله خوشامد گفته شود)، و رژیم حاکم می‌خواهد «ملیت و حاکمیت» را برجسته کند (تا مداخله نامشروع شود).


%% ================================================================
\section{اطاعت از مرجعیت: آزمایش میلگرام و کاربردهای سیاسی آن}
\label{sec:ch11-milgram}
%% ================================================================

\needspace{6\baselineskip}
\subsection{آزمایش اصلی}
\label{subsec:ch11-milgram-experiment}

در سال ۱۹۶۱، \thinker{استنلی میلگرام} (\lat{Stanley Milgram})، روان‌شناس دانشگاه ییل، آزمایش مشهور \concept{اطاعت از مرجعیت} (\lat{Obedience to Authority}) را انجام داد. شرکت‌کنندگان عادی به دستور یک «دانشمند» (بازیگر) حاضر شدند شوک‌های الکتریکی فزاینده (تا ۴۵۰ ولت — ظاهراً کشنده) به فرد دیگری (نیز بازیگر) وارد کنند. نتیجه‌ی شوکه‌کننده: ۶۵٪ شرکت‌کنندگان تا آخرین مرحله (۴۵۰ ولت) پیش رفتند.\footnote{\lr{Milgram, Stanley. \textit{Obedience to Authority: An Experimental View}. New York: Harper \& Row, 1974.}}

\begin{tcolorbox}[defbox, title={یافته‌ی کلیدی آزمایش میلگرام}]
انسان‌های عادی، در شرایطی که یک \concept{مرجعیت مشروع} (\lat{legitimate authority}) به آن‌ها دستور دهد، حاضرند اعمالی انجام دهند که در شرایط عادی هرگز انجام نمی‌دادند — حتی آسیب رساندن به انسان‌های بی‌گناه. عامل تعیین‌کننده نه شخصیت فرد، بلکه \textbf{ساختار موقعیت} و \textbf{مشروعیت ادراک‌شده‌ی مرجعیت} است.
\end{tcolorbox}

\subsection{کاربرد در نجات‌طلبی: مرجعیت خارجی}
\label{subsec:ch11-milgram-application}

آزمایش میلگرام برای فهم نجات‌طلبی از دو زاویه اهمیت دارد:

\textbf{زاویه‌ی اول: اطاعت از مرجعیت داخلی سرکوبگر.}
چرا مردم از رژیم‌های دیکتاتوری اطاعت می‌کنند، حتی وقتی اطاعت به زیان‌شان است؟ میلگرام نشان داد که اطاعت نه از «ضعف شخصیتی» بلکه از فشار موقعیتی ناشی می‌شود. همین مکانیزم، ناامیدی آموخته‌شده را تقویت می‌کند: مردم نه‌تنها توانایی مقاومت ندارند، بلکه \textit{حق} مقاومت را هم برای خود قائل نمی‌شوند.

\textbf{زاویه‌ی دوم: اطاعت از مرجعیت خارجی «نجات‌بخش».}
وقتی مرجعیت داخلی سرکوبگر مشروعیتش را از دست می‌دهد، مردم به‌دنبال \concept{مرجعیت جایگزین} می‌گردند. قدرت خارجی‌ای که وعده‌ی نجات می‌دهد، می‌تواند به مرجعیت جدید تبدیل شود — مرجعیتی که مردم حاضرند از آن اطاعت کنند و هزینه‌های مداخله‌اش را بپذیرند. \footnote{\lr{Blass, Thomas. \textit{The Man Who Shocked the World: The Life and Legacy of Stanley Milgram}. New York: Basic Books, 2004, pp. 261–279.}}

\begin{tcolorbox}[casebox, title={مطالعه‌ی موردی: اطاعت از مرجعیت در عراق ۲۰۰۳}]
پس از سقوط صدام، بخشی از جامعه‌ی عراق مرجعیت آمریکا و \histref{پل برمر}{مدیر اقتدار موقت ائتلاف} را پذیرفت — نه به دلیل عشق به آمریکا، بلکه به دلیل نیاز روان‌شناختی به مرجعیت. فرمان‌های برمر (انحلال ارتش، بعث‌زدایی) با مقاومت کمتری اجرا شدند

```latex
تا آنچه انتظار می‌رفت — زیرا بسیاری از عراقی‌ها، پس از دهه‌ها اطاعت از مرجعیت بعث، به‌سادگی مرجعیت جدید را جایگزین کردند.

اما مسئله اینجاست: مرجعیت خارجی — بر خلاف مرجعیت داخلی — فاقد \concept{پاسخگویی} (\lat{accountability}) به مردم محلی است. برمر به واشنگتن پاسخگو بود، نه به مردم بغداد. این شکاف پاسخگویی، زمینه‌ی سوءاستفاده از اطاعت را فراهم می‌آورد.\footnote{\lr{Diamond, Larry. \textit{Squandered Victory: The American Occupation and the Bungled Effort to Bring Democracy to Iraq}. New York: Times Books, 2005, pp. 29–46.}}
\end{tcolorbox}


%% ================================================================
\section{تفکر گروهی: وقتی جمع از فرد «احمق‌تر» می‌شود}
\label{sec:ch11-groupthink}
%% ================================================================

\needspace{6\baselineskip}
\subsection{نظریه‌ی جنیس}
\label{subsec:ch11-janis}

\thinker{ایروینگ جنیس} (\lat{Irving Janis}) در سال ۱۹۷۲ مفهوم \concept{تفکر گروهی} (\lat{Groupthink}) را مطرح کرد: پدیده‌ای که در آن فشار گروهی برای اجماع و هم‌نوایی، کیفیت تصمیم‌گیری را به‌شدت کاهش می‌دهد. در تفکر گروهی، اعضای گروه از بیان مخالفت خودداری می‌کنند، اطلاعات ناسازگار نادیده گرفته می‌شود، و توهم آسیب‌ناپذیری شکل می‌گیرد.\footnote{\lr{Janis, Irving L. \textit{Victims of Groupthink: A Psychological Study of Foreign-Policy Decisions and Fiascoes}. Boston: Houghton Mifflin, 1972.}}

\begin{tcolorbox}[defbox, title={تعریف: تفکر گروهی — هشت نشانه}]
\thinker{جنیس} هشت نشانه‌ی تفکر گروهی را شناسایی کرد:
\begin{enumerate}
\item \textbf{توهم آسیب‌ناپذیری:} «هیچ‌چیز نمی‌تواند به ما آسیب برساند»
\item \textbf{عقلانی‌سازی جمعی:} توجیه تصمیمات بد با استدلال‌های سطحی
\item \textbf{باور به اخلاق ذاتی گروه:} «ما طرف حق هستیم»
\item \textbf{کلیشه‌سازی درباره‌ی برون‌گروه:} «آن‌ها شرور/بی‌عقل/وحشی‌اند»
\item \textbf{فشار بر مخالفان:} کسانی که شک می‌کنند تحت فشار قرار می‌گیرند
\item \textbf{خودسانسوری:} اعضا از بیان تردیدهایشان خودداری می‌کنند
\item \textbf{توهم اجماع:} سکوت به معنای موافقت تلقی می‌شود
\item \textbf{محافظان ذهنی:} برخی اعضا خودبه‌خود اطلاعات ناسازگار را فیلتر می‌کنند
\end{enumerate}
\end{tcolorbox}

\subsection{تفکر گروهی در تصمیم به مداخله}
\label{subsec:ch11-groupthink-intervention}

جنیس خود نمونه‌های متعددی از تفکر گروهی در سیاست خارجی آمریکا را تحلیل کرد: حمله‌ی خلیج خوک‌ها (۱۹۶۱)، تشدید جنگ ویتنام، و واترگیت. پژوهشگران بعدی نشان داده‌اند که تصمیم به حمله به عراق در ۲۰۰۳ نمونه‌ی کلاسیک دیگری از تفکر گروهی بود.\footnote{\lr{Badie, Dina. "Groupthink, Iraq, and the War on Terror: Explaining US Policy Shift toward Iraq," \textit{Foreign Policy Analysis}, Vol. 6, No. 4, 2010, pp. 277–296.}}

\begin{tcolorbox}[casebox, title={مطالعه‌ی موردی: تفکر گروهی و تصمیم حمله به عراق ۲۰۰۳}]
\textbf{توهم آسیب‌ناپذیری:} پس از پیروزی سریع در افغانستان، اعتماد به نفس دولت بوش به اوج رسید.

\textbf{عقلانی‌سازی:} اطلاعات نادرست درباره‌ی سلاح‌های کشتار جمعی عقلانی‌سازی شدند. وقتی بازرسان سازمان ملل (\thinker{هانس بلیکس}، \lat{Hans Blix}) سلاحی نیافتند، این به معنای «پنهان‌کاری عراق» تفسیر شد، نه فقدان سلاح.\footnote{\lr{Blix, Hans. \textit{Disarming Iraq}. New York: Pantheon Books, 2004, pp. 194–219.}}

\textbf{باور اخلاقی:} نومحافظه‌کاران خود را «رهاکنندگان ملت‌ها» و ناقدان را «بزدل» یا «طرفدار صدام» می‌خواندند.

\textbf{کلیشه‌سازی:} صدام = هیتلر؛ عراق = آلمان نازی — قیاسی بسیار گمراه‌کننده.

\textbf{فشار بر مخالفان:} \thinker{اریک شینسکی} (\lat{Eric Shinseki})، فرمانده‌ی ارتش، که گفت اشغال عراق به «چند صد هزار» نیرو نیاز دارد، به حاشیه رانده شد. \thinker{جو ویلسون} (\lat{Joseph Wilson})، دیپلمات، که گزارش نادرستی اطلاعات هسته‌ای عراق را داد، مورد حمله قرار گرفت.\footnote{\lr{Ricks, Thomas E. \textit{Fiasco}, op. cit., pp. 96–97.}}

\textbf{خودسانسوری:} \thinker{کالین پاول} (\lat{Colin Powell})، وزیر خارجه، تردیدهای جدی درباره‌ی اطلاعات داشت اما در شورای امنیت به دفاع از آن‌ها پرداخت — تصمیمی که بعدها آن را «لکه‌ای بر کارنامه‌ام» نامید.\footnote{\lr{DeYoung, Karen. \textit{Soldier: The Life of Colin Powell}. New York: Alfred A. Knopf, 2006, pp. 401–435.}}
\end{tcolorbox}

\subsection{تفکر گروهی در ملت‌های نجات‌طلب}
\label{subsec:ch11-groupthink-nation}

تفکر گروهی فقط مختص مداخله‌گران نیست؛ در ملت‌های نجات‌طلب نیز عمل می‌کند. وقتی یک جامعه‌ی تحت فشار به اجماعی درباره‌ی «نیاز به نجات خارجی» می‌رسد، صداهای مخالف — کسانی که هشدار می‌دهند مداخله ممکن است بدتر باشد — سرکوب یا حاشیه‌رانی می‌شوند.

\thinker{کاس سانستین} (\lat{Cass Sunstein}) مفهوم \concept{قطب‌بندی گروهی} (\lat{Group Polarization}) را توسعه داده است: گروه‌هایی که هم‌فکرند، پس از بحث گروهی به مواضع \textit{افراطی‌تر} از میانگین اولیه‌ی اعضا می‌رسند.\footnote{\lr{Sunstein, Cass R. "The Law of Group Polarization," \textit{Journal of Political Philosophy}, Vol. 10, No. 2, 2002, pp. 175–195.}} در فضای شبکه‌های اجتماعی، قطب‌بندی تشدید می‌شود: الگوریتم‌ها محتوای تحریک‌آمیز را تقویت و صداهای میانه‌رو را حذف می‌کنند.\footnote{\lr{Bail, Christopher A., et al. "Exposure to Opposing Views on Social Media Can Increase Political Polarization," \textit{Proceedings of the National Academy of Sciences}, Vol. 115, No. 37, 2018, pp. 9216–9221.}}


%% ================================================================
\section{آزمایش زندان استنفورد و سوءاستفاده از قدرت}
\label{sec:ch11-stanford}
%% ================================================================

\needspace{6\baselineskip}
در سال ۱۹۷۱، \thinker{فیلیپ زیمباردو} (\lat{Philip Zimbardo}) آزمایش مشهور \concept{زندان استنفورد} (\lat{Stanford Prison Experiment}) را انجام داد. دانشجویان داوطلب به‌طور تصادفی به نقش «زندانبان» و «زندانی» تقسیم شدند. ظرف چند روز، «زندانبانان» رفتارهای سادیستیک و «زندانیان» رفتارهای منفعلانه و افسرده نشان دادند.\footnote{\lr{Zimbardo, Philip G. \textit{The Lucifer Effect: Understanding How Good People Turn Evil}. New York: Random House, 2007.}}

\begin{tcolorbox}[wavebox, title={درس زیمباردو برای مداخلات خارجی}]
زیمباردو بعد از رسوایی زندان \concept{ابوغریب} در عراق (۲۰۰۴) — جایی که سربازان آمریکایی زندانیان عراقی را شکنجه و تحقیر کردند — استدلال کرد که مسئله نه «سیب‌های گندیده» (افراد بد) بلکه «بشکه‌ی گندیده» (سیستم بد) است.\footnote{\lr{Zimbardo, Philip G. \textit{The Lucifer Effect}, op. cit., pp. 324–379.}}

\textbf{کاربرد در نجات‌طلبی:}
\begin{itemize}
\item هر مداخله‌ی خارجی یک \concept{نامتقارنی قدرت} ایجاد می‌کند: مداخله‌گر = زندانبان؛ ملت هدف = زندانی
\item حتی مداخله‌گران «خوش‌نیت» در این ساختار نامتقارن، مستعد سوءاستفاده از قدرت‌اند
\item ملت هدف نیز در موضع «زندانی» قرار می‌گیرد: ناامیدی آموخته‌شده تشدید می‌شود
\item ابوغریب استثنا نبود — تظاهر منطقی ساختار اشغال بود
\end{itemize}
\end{tcolorbox}


%% ================================================================
\section{نقشه‌ی ذهنی: مکانیزم‌های روان‌شناختی نجات‌طلبی}
\label{sec:ch11-mindmap}
%% ================================================================

\needspace{14\baselineskip}

\begin{figure}[htbp]
\centering
\begin{tikzpicture}[
    >=Stealth,
    mindmap,
    every node/.style={concept, font=\small, text=white, minimum size=1.5cm},
    root concept/.append style={concept color=PrimaryDeep, font=\normalsize\bfseries, minimum size=3cm},
    level 1 concept/.append style={sibling angle=60, level distance=4.5cm},
    level 2 concept/.append style={sibling angle=45, level distance=3cm, font=\tiny},
    scale=0.72, transform shape
]

\node[root concept] {\rl{نجات‌طلبی}}
    child[concept color=red!60!black] {
        node {\rl{ناامیدی}\\[2pt]\rl{آموخته‌شده}}
        child { node {\rl{نقص انگیزشی}} }
        child { node {\rl{نقص شناختی}} }
        child { node {\rl{نقص عاطفی}} }
    }
    child[concept color=orange!70!black] {
        node {\rl{آرزو-}\\[2pt]\rl{اندیشی}}
        child { node {\rl{سوگیری تأیید}} }
        child { node {\rl{سوگیری}\\[2pt]\rl{خوش‌بینی}} }
        child { node {\rl{سوگیری}\\[2pt]\rl{بازمانده}} }
    }
    child[concept color=green!50!black] {
        node {\rl{هویت}\\[2pt]\rl{قربانی}}
        child { node {\rl{تروما-ی}\\[2pt]\rl{انتخاب‌شده}} }
        child { node {\rl{مصونیت}\\[2pt]\rl{اخلاقی}} }
    }
    child[concept color=blue!60!black] {
        node {\rl{هویت}\\[2pt]\rl{اجتماعی}}
        child { node {\rl{مرز}\\[2pt]\rl{خودی/بیگانه}} }
        child { node {\rl{بازتعریف}\\[2pt]\rl{درون‌گروه}} }
    }
    child[concept color=purple!60!black] {
        node {\rl{تفکر}\\[2pt]\rl{گروهی}}
        child { node {\rl{توهم اجماع}} }
        child { node {\rl{قطب‌بندی}} }
    }
    child[concept color=brown!60!black] {
        node {\rl{اطاعت}\\[2pt]\rl{از مرجعیت}}
        child { node {\rl{مرجعیت}\\[2pt]\rl{داخلی}} }
        child { node {\rl{مرجعیت}\\[2pt]\rl{خارجی}} }
    };

\end{tikzpicture}
\caption{نقشه‌ی ذهنی مکانیزم‌های روان‌شناختی نجات‌طلبی}
\label{fig:ch11-mindmap}
\end{figure}


%% ================================================================
\section{شناخت اجتماعی و ادراک تهدید}
\label{sec:ch11-threat-perception}
%% ================================================================

\needspace{6\baselineskip}
\subsection{نظریه‌ی مدیریت وحشت}
\label{subsec:ch11-tmt}

\concept{نظریه‌ی مدیریت وحشت} (\lat{Terror Management Theory — TMT})، توسعه‌یافته توسط \thinker{جف گرینبرگ} (\lat{Jeff Greenberg})، \thinker{شلدون سولومون} (\lat{Sheldon Solomon})، و \thinker{تام پیژینسکی} (\lat{Tom Pyszczynski})، بر اساس آثار \thinker{ارنست بکر} (\lat{Ernest Becker}) بنا شده است. این نظریه ادعا می‌کند که آگاهی از مرگ، \concept{اضطراب وجودی} (\lat{Existential Anxiety}) عمیقی ایجاد می‌کند که انسان‌ها با تقویت \concept{جهان‌بینی فرهنگی} و \concept{عزت‌نفس} خود آن را مدیریت می‌کنند.\footnote{\lr{Greenberg, Jeff, Sheldon Solomon, and Tom Pyszczynski. "Terror Management Theory of Self-Esteem and Cultural Worldviews," in Mark P. Zanna (ed.), \textit{Advances in Experimental Social Psychology}, Vol. 29, San Diego: Academic Press, 1997, pp. 61–139.}}

\begin{tcolorbox}[defbox, title={مفهوم کلیدی: نظریه‌ی مدیریت وحشت و نجات‌طلبی}]
وقتی یک جامعه تحت تهدید وجودی قرار می‌گیرد (جنگ، سرکوب، فروپاشی اقتصادی)، \concept{برجستگی مرگ} (\lat{Mortality Salience}) افزایش می‌یابد. پاسخ روان‌شناختی شامل:
\begin{itemize}
\item \textbf{تقویت هویت گروهی:} پناه بردن به «ما» در برابر «آن‌ها»
\item \textbf{دفاع از جهان‌بینی:} تعصب نسبت به ارزش‌های خودی
\item \textbf{جستجوی رهبر کاریزماتیک:} نیاز به «نجات‌بخش» قوی و قاطع
\item \textbf{خشونت نسبت به برون‌گروه:} افزایش پرخاشگری نسبت به «دیگری»
\end{itemize}
آزمایش‌ها نشان داده‌اند که وقتی آزمودنی‌ها با یادآوری مرگ مواجه شوند، حمایت‌شان از رهبران اقتدارگرا و سیاست‌های تهاجمی افزایش می‌یابد.\footnote{\lr{Cohen, Florette, et al. "Fatal Attraction: The Effects of Mortality Salience on Evaluations of Charismatic, Task-Oriented, and Relationship-Oriented Leaders," \textit{Psychological Science}, Vol. 15, No. 12, 2004, pp. 846–851.}}
\end{tcolorbox}

کاربرد مستقیم در نجات‌طلبی: ملت‌هایی که تحت تهدید وجودی هستند — از رواندا ۱۹۹۴ تا سوریه ۲۰۱۱ — گرایش بیشتری به پذیرش مداخله‌ی خارجی نشان می‌دهند، زیرا اضطراب وجودی، نیاز به «نجات‌بخش» قوی و فوری را تشدید می‌کند. این همان مکانیزمی است که \concept{فرصت‌طلبان} — چه داخلی و چه خارجی — از آن بهره‌برداری می‌کنند.

\subsection{ادراک تهدید و اثر تهدید بر شناخت}
\label{subsec:ch11-threat-cognition}

\thinker{جنیفر لرنر} (\lat{Jennifer Lerner}) و همکارانش در پژوهش‌های پس از ۱۱ سپتامبر ۲۰۰۱ نشان دادند که \concept{ترس} و \concept{خشم} اثرات بسیار متفاوتی بر تصمیم‌گیری دارند: ترس، ریسک‌گریزی و محافظه‌کاری ایجاد می‌کند؛ خشم، ریسک‌پذیری و تمایل به اقدام تهاجمی.\footnote{\lr{Lerner, Jennifer S., et al. "Effects of Fear and Anger on Perceived Risks of Terrorism: A National Field Experiment," \textit{Psychological Science}, Vol. 14, No. 2, 2003, pp. 144–150.}}

\begin{tcolorbox}[wavebox, title={تحلیل: ترس، خشم، و نجات‌طلبی}]
\textbf{ملت‌های ترسیده} گرایش دارند هر وعده‌ی «امنیت» را بپذیرند — حتی اگر از سوی قدرت خارجی باشد. ترس عقلانیت تحلیلی را کاهش و نیاز به «پناهگاه» را افزایش می‌دهد.

\textbf{ملت‌های خشمگین} گرایش دارند خواستار اقدام فوری شوند — حتی اگر اقدام نسنجیده باشد. خشم، ارزیابی ریسک را کاهش و تمایل به «انتقام» یا «عمل قاطع» را افزایش می‌دهد.

\textbf{مداخله‌گران} از هر دو احساس بهره‌برداری می‌کنند: ترس مردم هدف از رژیم حاکم + خشم آن‌ها نسبت به سرکوب = آمادگی برای پذیرش «هر تغییری» — حتی اشغال خارجی.
\end{tcolorbox}


%% ================================================================
\section{نقش رسانه و تبلیغات در شکل‌دهی روان‌شناسی نجات‌طلبی}
\label{sec:ch11-propaganda}
%% ================================================================

\needspace{6\baselineskip}
\subsection{مدل تبلیغاتی هرمان و چامسکی}
\label{subsec:ch11-propaganda-model}

\thinker{ادوارد هرمان} (\lat{Edward S. Herman}) و \thinker{نوام چامسکی} (\lat{Noam Chomsky}) در کتاب \textit{تولید رضایت} مدل \concept{تبلیغات} (\lat{Propaganda Model}) را ارائه کردند. بر اساس این مدل، رسانه‌های جریان اصلی — حتی در جوامع دموکراتیک — از پنج «فیلتر» عبور می‌کنند که محتوا را به نفع نخبگان قدرت شکل می‌دهند: (۱) مالکیت، (۲) تبلیغات به‌عنوان منبع درآمد، (۳) وابستگی به منابع خبری رسمی، (۴) «ضربات انضباطی» (\lat{flak})، و (۵) ایدئولوژی ضدکمونیسم (یا در دوره‌ی معاصر: ضدتروریسم).\footnote{\lr{Herman, Edward S. and Noam Chomsky. \textit{Manufacturing Consent: The Political Economy of the Mass Media}. New York: Pantheon Books, 1988; updated ed., 2002.}}

\subsection{تبلیغات جنگی و ساختن «قربانی شایسته»}
\label{subsec:ch11-worthy-victims}

هرمان و چامسکی مفهوم \concept{قربانی شایسته} (\lat{Worthy Victim}) و \concept{قربانی ناشایسته} (\lat{Unworthy Victim}) را مطرح کردند: رسانه‌ها قربانیان «دشمنان» غرب را برجسته و قربانیان خود غرب یا متحدانش را کم‌اهمیت جلوه می‌دهند. این انتخابی بودن رسانه‌ای مستقیماً بر روان‌شناسی نجات‌طلبی تأثیر دارد.

\begin{table}[htbp]
\centering
\caption{قربانیان شایسته و ناشایسته در پوشش رسانه‌ای}
\label{tab:ch11-worthy-victims}
\begin{adjustbox}{max width=\linewidth}
\begin{tabularx}{\linewidth}{
    >{\raggedleft\arraybackslash}p{3cm}
    >{\raggedleft\arraybackslash}p{3cm}
    >{\raggedleft\arraybackslash}X
}
\toprule
\rowcolor{PrimaryDeep}
\textcolor{white}{\textbf{قربانی «شایسته»}} &
\textcolor{white}{\textbf{قربانی «ناشایسته»}} &
\textcolor{white}{\textbf{توضیح}} \\
\midrule
قربانیان صدام حسین & قربانیان تحریم‌های عراق & کشتار صدام پوشش وسیع؛ مرگ کودکان بر اثر تحریم کم‌پوشش \\
\rowcolor{NeutralLight}
قربانیان بمباران حلب (روسیه) & قربانیان بمباران موصل (آمریکا) & بمباران روسیه «جنایت جنگی»؛ بمباران آمریکا «دقیق و هدفمند» \\
قربانیان سرکوب قذافی & قربانیان هرج‌ومرج پس از قذافی & پیش از مداخله: پوشش وسیع؛ پس از مداخله: فراموشی \\
\rowcolor{NeutralLight}
قربانیان حمله‌ی روسیه به اوکراین & قربانیان جنگ یمن (عربستان) & اوکراین: پوشش ۲۴ ساعته؛ یمن: «جنگ فراموش‌شده» \\
\bottomrule
\end{tabularx}
\end{adjustbox}
\end{table}

\subsection{رسانه‌های فارسی‌زبان خارجی و نجات‌طلبی ایرانی}
\label{subsec:ch11-persian-media}

یکی از نمونه‌های مشخص تأثیر رسانه بر روان‌شناسی نجات‌طلبی، نقش رسانه‌های فارسی‌زبان خارجی (بی‌بی‌سی فارسی، صدای آمریکا، رادیو فردا، ایران اینترنشنال) در شکل‌دهی به ادراک بخشی از جامعه‌ی ایرانی از مداخله‌ی خارجی است. بدون قضاوت ارزشی، می‌توان مکانیزم‌های روان‌شناختی زیر را شناسایی کرد:

\begin{enumerate}
\item \textbf{برجسته‌سازی نارضایتی:} تأکید بر مشکلات و نارسایی‌ها — هم‌راستا با سوگیری منفی‌گرایی (\lat{Negativity Bias})\footnote{\lr{Rozin, Paul and Edward B. Royzman. "Negativity Bias, Negativity Dominance, and Contagion," \textit{Personality and Social Psychology Review}, Vol. 5, No. 4, 2001, pp. 296–320.}}
\item \textbf{الگوسازی از «موفقیت‌های» خارجی:} معرفی مدل‌های غربی به‌عنوان تنها مسیر پیشرفت
\item \textbf{حاشیه‌رانی صداهای میانه‌رو:} تقویت صداهای افراطی (رژیم‌چنج) و تضعیف صداهای اصلاح‌طلب
\item \textbf{ایجاد حس فوریت:} «فرصت تاریخی»، «الان یا هرگز» — مشابه تکنیک‌های فروش\footnote{\lr{Cialdini, Robert B. \textit{Influence: The Psychology of Persuasion}. Rev. ed., New York: HarperBusiness, 2006.}}
\end{enumerate}


%% ================================================================
\section{مقاومت در برابر نجات‌طلبی: عاملیت و تاب‌آوری}
\label{sec:ch11-resilience}
%% ================================================================

\needspace{6\baselineskip}
تصویر یکسره تاریکی که تا اینجا ترسیم شد — ناامیدی، سوگیری، تروما، اطاعت — ممکن است این تصور نادرست را ایجاد کند که ملت‌ها محکوم به نجات‌طلبی‌اند. اما روان‌شناسی مثبت و تحقیقات تاب‌آوری نشان می‌دهند که مقاومت در برابر این مکانیزم‌ها ممکن است.

\subsection{مفهوم عاملیت جمعی}
\label{subsec:ch11-agency}

\concept{عاملیت جمعی} (\lat{Collective Agency}) — باور مشترک اعضای یک گروه به توانایی‌شان در ایجاد تغییر — نقطه‌ی مقابل ناامیدی آموخته‌شده است. \thinker{آلبرت بندورا} (\lat{Albert Bandura})، روان‌شناس شناختی-اجتماعی، مفهوم \concept{خودکارآمدی جمعی} (\lat{Collective Efficacy}) را مطرح کرده است: باور گروه به اینکه می‌تواند از طریق اقدام مشترک به اهدافش دست یابد.\footnote{\lr{Bandura, Albert. "Exercise of Human Agency Through Collective Efficacy," \textit{Current Directions in Psychological Science}, Vol. 9, No. 3, 2000, pp. 75–78.}}

\begin{tcolorbox}[successbox, title={عوامل تقویت عاملیت جمعی (در مقابل نجات‌طلبی)}]
\begin{enumerate}
\item \textbf{روایت‌های موفقیت داخلی:} یادآوری نمونه‌هایی که ملت بدون کمک خارجی موفق شده (مثلاً: مشروطه‌ی ایران ۱۹۰۶، مقاومت ویتنام، جنبش ضد آپارتاید)
\item \textbf{رهبری داخلی کاریزماتیک:} وجود رهبران بومی که مشروعیت و اعتماد دارند (ماندلا، گاندی)
\item \textbf{شبکه‌های اجتماعی مستقل:} نهادهای مدنی، سندیکاها، انجمن‌های صنفی، مساجد/کلیساها
\item \textbf{تجربه‌های موفق مقاومت مدنی:} \thinker{اریکا چنووت} (\lat{Erica Chenoweth}) نشان داده است که مقاومت مدنی غیرخشونت‌آمیز دو برابر بیشتر از مقاومت مسلحانه به موفقیت می‌رسد\footnote{\lr{Chenoweth, Erica and Maria J. Stephan. \textit{Why Civil Resistance Works: The Strategic Logic of Nonviolent Conflict}. New York: Columbia University Press, 2011.}}
\item \textbf{آموزش تفکر انتقادی:} توانایی شناسایی سوگیری‌ها، تبلیغات، و اطلاعات جعلی
\item \textbf{حافظه‌ی تاریخی متوازن:} یادآوری هم پیروزی‌ها و هم شکست‌ها — نه فقط تروما
\end{enumerate}
\end{tcolorbox}

\subsection{مقاومت اوکراین: نمونه‌ی عاملیت جمعی}
\label{subsec:ch11-ukraine-agency}

همان‌گونه که در فصل~\ref{ch:contemporary} بحث شد، مقاومت اوکراین در برابر تهاجم روسیه (۲۰۲۲) نمونه‌ی بارز عاملیت جمعی بود. اوکراین از غرب کمک گرفت (تسلیحات، اطلاعات، تحریم‌ها)، اما عامل اصلی مقاومت، اراده‌ی ملی و خودکارآمدی جمعی بود. تفاوت کلیدی اوکراین با عراق و افغانستان: مردم اوکراین خود جنگیدند و از غرب فقط ابزار خواستند — نه «نجات‌بخش». این تفاوت بنیادین میان \concept{کمک‌خواهی} و \concept{نجات‌طلبی} است.\footnote{\lr{Marples, David R. "Understanding Ukraine's Euromaidan and Revolution of Dignity," in \textit{The Conflict in Ukraine: What Everyone Needs to Know}. Oxford: Oxford University Press, 2023, pp. 89–112.}}


%% ================================================================
\section{نمودار: چرخه‌ی روان‌شناختی نجات‌طلبی}
\label{sec:ch11-cycle-tikz}
%% ================================================================

\needspace{14\baselineskip}

\begin{figure}[htbp]
\centering
\begin{tikzpicture}[
    >=Stealth,
    node distance=3cm,
    phase/.style={rectangle, draw=PrimaryDeep, fill=PrimaryDeep!10, rounded corners=8pt,
        text width=3.2cm, minimum height=1.5cm, align=center, font=\small},
    break/.style={rectangle, draw=green!60!black, fill=green!10, rounded corners=8pt,
        text width=3cm, minimum height=1cm, align=center, font=\small},
    arrow/.style={->, very thick, PrimaryDeep},
    breakarrow/.style={->, very thick, green!60!black, dashed},
    scale=0.82, transform shape
]

% چرخه‌ی اصلی (دایره‌ای)
\node[phase] (oppression) at (90:5cm) {\rl{سرکوب / بحران}\\[3pt]\rl{(محرک اولیه)}};
\node[phase] (helpless) at (30:5cm) {\rl{ناامیدی}\\[3pt]\rl{آموخته‌شده}};
\node[phase] (wishful) at (330:5cm) {\rl{آرزواندیشی}\\[3pt]\rl{و سوگیری‌ها}};
\node[phase] (savior) at (270:5cm) {\rl{جستجوی}\\[3pt]\rl{نجات‌بخش خارجی}};
\node[phase] (intervention) at (210:5cm) {\rl{مداخله‌ی خارجی}};
\node[phase] (failure) at (150:5cm) {\rl{شکست / پیامدهای}\\[3pt]\rl{ناخواسته}};

% فلش‌های چرخه
\draw[arrow] (oppression) -- (helpless);
\draw[arrow] (helpless) -- (wishful);
\draw[arrow] (wishful) -- (savior);
\draw[arrow] (savior) -- (intervention);
\draw[arrow] (intervention) -- (failure);
\draw[arrow] (failure) -- (oppression) node[midway, left, font=\tiny, text=red!70!black] {\rl{تکرار چرخه}};

% نقطه‌ی خروج
\node[break] (agency) at (0:0) {\rl{عاملیت جمعی}\\[3pt]\rl{(نقطه‌ی خروج)}};

\draw[breakarrow] (helpless) -- (agency) node[midway, above, font=\tiny, text=green!60!black] {\rl{شکستن ناامیدی}};
\draw[breakarrow] (wishful) -- (agency) node[midway, right, font=\tiny, text=green!60!black] {\rl{تفکر انتقادی}};
\draw[breakarrow] (savior) -- (agency) node[midway, below, font=\tiny, text=green!60!black] {\rl{اتکا به خود}};

% عنوان
\node[font=\normalsize\bfseries, PrimaryDeep] at (0, 7) {\rl{چرخه‌ی روان‌شناختی نجات‌طلبی و نقطه‌ی خروج}};

\end{tikzpicture}
\caption{چرخه‌ی روان‌شناختی نجات‌طلبی: از سرکوب تا شکست و بازگشت}
\label{fig:ch11-cycle}
\end{figure}


%% ================================================================
\section{جمع‌بندی فصل}
\label{sec:ch11-conclusion}
%% ================================================================

\needspace{6\baselineskip}
این فصل نشان داد که نجات‌طلبی — یعنی امیدبستن به نیروی خارجی برای رهایی از وضعیت نامطلوب — فقط یک «انتخاب سیاسی» نیست، بلکه ریشه‌های عمیقی در \concept{روان‌شناسی فردی و اجتماعی} دارد. مکانیزم‌هایی چون ناامیدی آموخته‌شده، آرزواندیشی، سوگیری‌های شناختی، هویت قربانی، تروماهای جمعی، بازتعریف مرز خودی/بیگانه، تفکر گروهی، و اطاعت از مرجعیت، همه در کنار هم زمینه‌ای روان‌شناختی فراهم می‌کنند که ملت‌ها را به سوی نجات‌طلبی سوق می‌دهد.

\begin{tcolorbox}[wavebox, title={یافته‌های کلیدی فصل ۱۱}]
\begin{enumerate}
\item \textbf{ناامیدی آموخته‌شده‌ی جمعی:} ملت‌هایی که تحت سرکوب طولانی‌مدت بوده‌اند، ممکن است باور به توانایی تغییر از درون را از دست بدهند. این ناامیدی، آن‌ها را مستعد پذیرش «نجات‌بخش خارجی» می‌سازد — حتی با شواهد تاریخی مبنی بر شکست مداخلات قبلی.

\item \textbf{آرزواندیشی:} سوگیری‌های شناختی — از جمله سوگیری تأیید، خوش‌بینی، و بازمانده — ارزیابی عقلانی هزینه-فایده‌ی مداخله را مختل می‌کنند. مردم آنچه را \textit{می‌خواهند} واقعیت باشد، با آنچه \textit{هست} اشتباه می‌گیرند.

\item \textbf{تروماهای جمعی:} تروماهای انتخاب‌شده (ولکان) هویت قربانی‌گری را بازتولید می‌کنند و فضای روانی را برای پذیرش مداخله آماده می‌سازند. تکرار اجباری فرویدی در سطح جمعی، ملت‌ها را در چرخه‌ی شکست نگه می‌دارد.

\item \textbf{بازتعریف هویت:} نظریه‌ی هویت اجتماعی نشان می‌دهد که مرز خودی/بیگانه سیّال است و مداخله‌گران از این سیّالیت برای مشروعیت‌بخشی به مداخله بهره‌برداری می‌کنند.

\item \textbf{تفکر گروهی:} هم در تصمیم‌گیری مداخله‌گران (عراق ۲۰۰۳) و هم در ملت‌های نجات‌طلب، تفکر گروهی کیفیت تحلیل را کاهش و ریسک شکست را افزایش می‌دهد.

\item \textbf{سوءاستفاده‌پذیری:} مکانیزم‌های روان‌شناختی فوق قابل شناسایی و بنابراین قابل بهره‌برداری هستند: مداخله‌گران (و فرصت‌طلبان داخلی) از ناامیدی، ترس، خشم، و آرزواندیشی ملت‌ها به‌عنوان \concept{ابزار مشروعیت‌سازی} استفاده می‌کنند.

\item \textbf{نقطه‌ی خروج:} عاملیت جمعی، خودکارآمدی، تفکر انتقادی، مقاومت مدنی غیرخشونت‌آمیز، و روایت‌های متوازن تاریخی می‌توانند چرخه‌ی نجات‌طلبی را بشکنند. تفاوت بنیادین میان \concept{کمک‌خواهی} (درخواست ابزار از خارج ضمن حفظ عاملیت داخلی — مدل اوکراین) و \concept{نجات‌طلبی} (واگذاری عاملیت به خارجی — مدل عراق) کلید فهم این تمایز است.

\item \textbf{پل به فصل بعد:} فصل ۱۲ ابعاد جامعه‌شناختی و اقتصادی نجات‌طلبی را بررسی خواهد کرد — نشان خواهد داد که مکانیزم‌های روان‌شناختی بررسی‌شده در این فصل در بستری اجتماعی-اقتصادی عمل می‌کنند و بدون فهم آن بستر، تحلیل ناقص خواهد ماند.
\end{enumerate}
\end{tcolorbox}

\begin{tcolorbox}[quotebox, title={سخن پایانی فصل}]
«هیچ‌کس نمی‌تواند شما را بدون رضایت خودتان احساس حقارت بدهد.»\\
— \thinker{النور روزولت} (\lat{Eleanor Roosevelt})\footnote{\lr{Roosevelt, Eleanor. \textit{This Is My Story}. New York: Harper \& Brothers, 1937, p. 192.}}\\[6pt]
«و هیچ‌کس نمی‌تواند ملتی را "نجات" دهد که خود نجات خویش را نمی‌خواهد — یا نمی‌داند.»\\
— نویسنده
\end{tcolorbox}

\cleardoublepage